\uuid{Z2i8}
\exo7id{1624}
\titre{exo7 1624}
\auteur{barraud}
\organisation{exo7}
\datecreate{2003-09-01}
\isIndication{false}
\isCorrection{false}
\chapitre{Réduction d'endomorphisme, polynôme annulateur}
\sousChapitre{Diagonalisation}
\cpge{mpsi-mp2i.applications-lineaires.a, pcsi.applications-lineaires.a, ptsi.applications-lineaires.a, mpsi-mp2i.applications-lineaires.b, pcsi.applications-lineaires.b, ptsi.applications-lineaires.b, mpsi-mp2i.matrices-applications-lineaires.a, pcsi.matrices-applications-lineaires.a, ptsi.matrices-applications-lineaires.a, bcpst1.applications-lineaires-matrices, bcpst1.matrices, mp-mpi.reduction-endomorphismes.b, pc.reduction.a, psi.reduction-endomorphismes.a, pt.reduction-endomorphismes.a, mp-mpi.reduction-endomorphismes.d, pc.reduction.c, psi.reduction-endomorphismes.c, pt.reduction-endomorphismes.c, bcpst2.valeurs-propres-vecteurs-propres.a, bcpst2.valeurs-propres-vecteurs-propres.b}
\module{Algèbre}
\niveau{L2}
\difficulte{}

\contenu{
\question{
Soit $u$ l'application suivante :
$$
  u:
  \begin{array}{rcl}
    \R_{2}[X] & \rightarrow  & \R_{2}[X]  \\[2ex]
    P         & \mapsto & (2X+1)P-(X^{2}-1)P'
  \end{array}
$$
Montrer que $u$ est bien d\'{e}finie et lin\'{e}aire. 
D\'{e}terminer les valeurs propres de $u$, et, si c'est possible, diagonaliser $u$.
}
}
